\scanpage{245}
\mspage{191}

\margin{$\Vert$}
On a \deleted{donc} vu au moins ceci : \unsure{Si}\textsuperscript{*} $X$ est cofibrée donc $W$-propre sur $S$, il ne s'ensuit nullement que $\unsure{\Phi|\mathscr{D}}(X)$ soit \deleted{\unsure{cofibrée}} \inserted{même $W$-propre} sur $S^{*}$, ni que ce soit une $W$-fibrée.
% slanted
\margin{voir résultat plus précis \unsure{page suivante|pages suivantes} !}

Mais que se passe-t-il si $X\in\underline{\mathrm{Prop}}_{\unsure{\circ|0}W}S$, par exemple s'il \deleted{\unsure{provient}} vient de $\underline{\mathrm{Cofib\,sc}}_{\circ}S$ ? Alors
\deleted{\unsure{pour} \illegible{1}} \deleted{\diagram{\begin{tikzcd}[ampersand replacement=\&,column sep=small,row sep=tiny]
\&\unsure{a}\\
\unsure{0}\arrow[ur]\arrow[dr]\&\\
\&b
\end{tikzcd}}} $X$ est dans $\underline{\mathrm{Fib}}_W S$, donc
on sait déjà \unsure{que} $\Phi(X)$ est dans $\underline{\mathrm{Fib\,sc}}_{\circ W}S$, donc ce n'est pas par là qu'on fera un contre-exemple. Mais considérons un un $X$ cofibrée sur

\diagram{\begin{tikzcd}[ampersand replacement=\&,column sep=small,row sep=tiny]
\&a\\
S=0\arrow[ur]\arrow[dr]\&\\
\&b\quad,
\end{tikzcd}}

donnée par le diagramme

\diagram{\begin{tikzcd}[ampersand replacement=\&,column sep=small,row sep=tiny]
\&X_{\unsure{a|0}}\\
X_0\arrow[ur,"u"]\arrow[dr,"\unsure{v|u}"]\&\\
\&X_b
\end{tikzcd}}

dans $\mathrm{Cat}$, en supposant cette fois $u,v\in W$ ($X\in\underline{\mathrm{Cofib\,sc}}_{\circ W}S$). \deleted{\unsure{D'où $X$ est telle que}} Alors $\Phi(X)$ est la catégorie fibrée scindée sur $S$ définie par le diagramme contravariant sur $S$

\diagram{\begin{tikzcd}[ampersand replacement=\&,column sep=small,row sep=tiny]
\&X_a\arrow[dl,hook,"i"']\\
X\&\\
\&X_b\arrow[ul,hook,"j"]
\end{tikzcd}}\quad,

où $i,j$ sont les inclusions (qui sont des imm. fermées). Dire que $\unsure{\mathscr{X}|\mathfrak{X}}=\Phi(X)$ est propre sur $S$ signifie aussi que ses restrictions au dessus des \deleted{\unsure{les}} sous-catégories ouvertes $\{0,a\}$ et $\{0,b\}$ de $S$ le sont (car toute flèche $\Delta^1\to S$ se factorise par l'une ou l'autre). Cela
% archive sheet number: 244
